%=============================================================================
% CP2_Multiplicities_Definitive.tex
% Definitive resolution of eigenvalue multiplicities for the Laplacian on CP^2
% Date: 2026-03-22
%=============================================================================

\documentclass[11pt]{article}
\usepackage{amsmath,amssymb,amsthm}
\usepackage{booktabs}
\usepackage{geometry}
\geometry{margin=1in}

\newtheorem{theorem}{Theorem}
\newtheorem{proposition}{Proposition}
\newtheorem{corollary}{Corollary}
\theoremstyle{remark}
\newtheorem{remark}{Remark}

\title{Definitive Eigenvalue Multiplicities for the Scalar Laplacian on
$\mathbb{C}P^2$: Resolution of the $(k{+}1)^2$ vs.\ $(k{+}1)^3$ Discrepancy}
\author{DFD Cross-Reference Audit}
\date{22 March 2026}

\begin{document}
\maketitle

%=============================================================================
\section{Executive Summary}
%=============================================================================

\begin{center}
\fbox{\parbox{0.9\textwidth}{%
\textbf{VERDICT:} The correct multiplicity for the scalar Laplacian on
$\mathbb{C}P^2$ is
\[
\boxed{m_k = (k+1)^3}\,.
\]
The numerical code in \texttt{numerical\_a4.py} and
\texttt{a4\_Gilkey\_Computation.tex} that uses $(k+1)^2$ is computing the
spectrum of the \textbf{3-sphere $S^3$}, not $\mathbb{C}P^2$.  Both manifolds
share the same eigenvalues $\lambda_k = k(k+2)$, but their multiplicities
differ.  Agent~3's claim of $(k+1)^3$ is correct for $\mathbb{C}P^2$.
}}
\end{center}

%=============================================================================
\section{The Problem}
%=============================================================================

Two different multiplicity formulas appear in the DFD codebase for the scalar
Laplacian $\Delta$ on $\mathbb{C}P^2$:
\begin{enumerate}
\item \textbf{Agent~3 / Toeplitz\_Heat\_Kernel.tex:} $m_k = (k+1)^3$
\item \textbf{numerical\_a4.py / a4\_Gilkey\_Computation.tex:} $m_k = (k+1)^2$
\end{enumerate}
Both use the same eigenvalues $\lambda_k = k(k+2)$ for $k = 0, 1, 2, \ldots$

The question is: which is the correct multiplicity for $\mathbb{C}P^2$
(real dimension~4)?

%=============================================================================
\section{Resolution: Three Independent Derivations}
%=============================================================================

\subsection{Derivation 1: Representation Theory of $\mathrm{SU}(3)$}

$\mathbb{C}P^2$ is the symmetric space $\mathrm{SU}(3)/\mathrm{U}(2)$.
By the Peter--Weyl theorem, the $L^2$ functions on $G/K$ decompose into
spherical (class-1) representations of $G$ with respect to $K$.

\paragraph{Step 1: Identify the spherical representations.}
For the symmetric space $\mathrm{SU}(n{+}1)/\mathrm{U}(n)$, the spherical
representations of $\mathrm{SU}(n{+}1)$ are precisely those with Dynkin labels
$(k, 0, \ldots, 0, k)$ for $k = 0, 1, 2, \ldots$ (the ``symmetric square''
representations).

For $\mathrm{SU}(3)$, which has rank~2, the Dynkin labels are $(a,b)$, so the
spherical representations are $(k,k)$ for $k \geq 0$.

\paragraph{Step 2: Compute the Casimir eigenvalue.}
The quadratic Casimir operator $C_2$ of $\mathrm{SU}(3)$ acts on the
irreducible representation $(a,b)$ with eigenvalue
\begin{equation}
C_2(a,b) = \frac{a^2 + ab + b^2}{3} + a + b\,.
\end{equation}
For the spherical representation $(k,k)$:
\begin{equation}
C_2(k,k) = \frac{k^2 + k^2 + k^2}{3} + 2k = k^2 + 2k = k(k+2)\,.
\end{equation}
This matches the eigenvalue $\lambda_k = k(k+2)$.

(Note: with the standard Fubini--Study normalization giving
$\lambda_k = 4k(k+n)$, one has $n=2$, so $4k(k+2)$.  The factor of~4
depends on the curvature normalization.  We use $\lambda_k = k(k+2)$
throughout, corresponding to Ricci curvature $\mathrm{Ric} = 2g$.)

\paragraph{Step 3: Compute the dimension (multiplicity).}
The Weyl dimension formula for $\mathrm{SU}(3)$ gives:
\begin{equation}
\dim(a,b) = \frac{(a+1)(b+1)(a+b+2)}{2}\,.
\end{equation}
For $(k,k)$:
\begin{equation}
\dim(k,k) = \frac{(k+1)(k+1)(2k+2)}{2}
           = \frac{(k+1)^2 \cdot 2(k+1)}{2}
           = (k+1)^3\,.
\end{equation}

\begin{proposition}
The multiplicity of the eigenvalue $\lambda_k = k(k+2)$ of the scalar
Laplacian on $\mathbb{C}P^2$ is $m_k = (k+1)^3$.
\end{proposition}

\subsection{Derivation 2: Bi-Homogeneous Harmonic Polynomials}

The eigenfunctions of the Laplacian on $\mathbb{C}P^n$ can be realized as the
restrictions to $\mathbb{C}P^n$ of bi-homogeneous harmonic polynomials
$P(z, \bar{z})$ on $\mathbb{C}^{n+1}$ of bi-degree $(k,k)$---i.e.,
homogeneous of degree~$k$ in the holomorphic variables $z_0, \ldots, z_n$
and degree~$k$ in the anti-holomorphic variables
$\bar{z}_0, \ldots, \bar{z}_n$.

The condition of equal bi-degree $(k,k)$ is required so that $P$ descends
to a well-defined function on $\mathbb{C}P^n = \mathbb{C}^{n+1}\!/\mathbb{C}^*$
(the phase cancels: $P(\lambda z, \bar{\lambda}\bar{z})
= |\lambda|^{2k} P(z,\bar{z})$).

\paragraph{Counting.}
The space of all bi-homogeneous polynomials of bi-degree $(k,k)$ on
$\mathbb{C}^{n+1}$ has dimension
\begin{equation}
\binom{k+n}{n}^{\!2}\,,
\end{equation}
since we independently choose a monomial of degree~$k$ in $(n{+}1)$
holomorphic variables and one in $(n{+}1)$ anti-holomorphic variables.

The harmonic subspace (annihilated by $\Delta_{\mathbb{C}^{n+1}}$)
has dimension
\begin{equation}
m_k = \binom{k+n}{n}^{\!2} - \binom{k+n-1}{n}^{\!2}\,.
\end{equation}

For $\mathbb{C}P^2$ ($n = 2$):
\begin{align}
m_k &= \binom{k+2}{2}^{\!2} - \binom{k+1}{2}^{\!2} \\
    &= \left[\frac{(k+1)(k+2)}{2}\right]^{\!2}
     - \left[\frac{k(k+1)}{2}\right]^{\!2} \\
    &= \frac{(k+1)^2}{4}\bigl[(k+2)^2 - k^2\bigr] \\
    &= \frac{(k+1)^2}{4}\cdot 4(k+1) \\
    &= (k+1)^3\,.
\end{align}

This confirms the result independently.

\subsection{Derivation 3: General Formula and the $\mathbb{C}P^n$ Family}

The general multiplicity formula for $\mathbb{C}P^n$ admits the compact form:
\begin{equation}
\boxed{m_k(\mathbb{C}P^n) = \frac{2k+n}{n}\,\binom{k+n-1}{n-1}^{\!2}}
\end{equation}

\noindent
This is equivalent to $\binom{k+n}{n}^2 - \binom{k+n-1}{n}^2$ by the
identity:
\begin{equation}
\binom{k+n}{n}^{\!2} - \binom{k+n-1}{n}^{\!2}
= \frac{2k+n}{n}\,\binom{k+n-1}{n-1}^{\!2}\,.
\end{equation}

\paragraph{Consistency checks.}
\begin{center}
\begin{tabular}{clll}
\toprule
$n$ & Manifold & $m_k$ formula & Simplification \\
\midrule
$1$ & $\mathbb{C}P^1 = S^2$ & $(2k+1)\cdot 1^2 = 2k+1$
    & Standard spherical harmonics \\
$2$ & $\mathbb{C}P^2$ & $\frac{2k+2}{2}(k+1)^2 = (k+1)^3$
    & Confirmed above \\
$3$ & $\mathbb{C}P^3$ & $\frac{2k+3}{3}\binom{k+2}{2}^2$
    & Not a perfect power \\
\bottomrule
\end{tabular}
\end{center}

The $n=1$ case is an essential sanity check: $\mathbb{C}P^1 \cong S^2$, and
the formula gives $m_k = 2k+1$, which is the well-known multiplicity for
eigenvalue $\lambda_k = k(k+1)$ of the Laplacian on $S^2$.

%=============================================================================
\section{Weyl's Law Verification}
%=============================================================================

Weyl's law for the eigenvalue counting function on a compact Riemannian
$d$-manifold states:
\begin{equation}
N(\Lambda) = \#\{\lambda_k \leq \Lambda\}
\sim \frac{\operatorname{Vol}(M)}{(4\pi)^{d/2}\,\Gamma(d/2+1)}\,
\Lambda^{d/2}
\quad\text{as }\Lambda\to\infty\,.
\end{equation}

For $\mathbb{C}P^2$, $d = 4$, so $N(\Lambda) \sim c\,\Lambda^2$.

\paragraph{With $(k+1)^3$ multiplicities (correct, $\mathbb{C}P^2$):}
Using $\lambda_k = k(k+2) \sim k^2$ for large $k$, we have
$k_{\max} \sim \Lambda^{1/2}$ and
\begin{equation}
N(\Lambda) = \sum_{k=0}^{k_{\max}} (k+1)^3
\sim \frac{k_{\max}^4}{4}
\sim \frac{\Lambda^2}{4}\,,
\end{equation}
which scales as $\Lambda^2$---\textbf{correct for dimension~4}.

\paragraph{With $(k+1)^2$ multiplicities (incorrect for $\mathbb{C}P^2$,
correct for $S^3$):}
\begin{equation}
N(\Lambda) = \sum_{k=0}^{k_{\max}} (k+1)^2
\sim \frac{k_{\max}^3}{3}
\sim \frac{\Lambda^{3/2}}{3}\,,
\end{equation}
which scales as $\Lambda^{3/2}$---correct for \textbf{dimension~3} ($S^3$),
not dimension~4.

\paragraph{Numerical verification.}
\begin{center}
\begin{tabular}{rrrcrc}
\toprule
$\Lambda$ & $N$ with $(k+1)^2$ & $N/\Lambda^{3/2}$
          & $N$ with $(k+1)^3$ & $N/\Lambda^2$ & Weyl dim \\
\midrule
   100 &       385 & 0.385 &      3,025 & 0.3025 & \\
 1,000 &    10,416 & 0.329 &    246,016 & 0.2460 & \\
10,000 &   338,350 & 0.338 & 25,502,500 & 0.2550 & \\
\bottomrule
\end{tabular}
\end{center}

The ratios $N/\Lambda^{3/2}$ stabilize for $(k+1)^2$, confirming 3D behavior.
The ratios $N/\Lambda^2$ stabilize for $(k+1)^3$, confirming 4D behavior.
Only $(k+1)^3$ is consistent with the real dimension of $\mathbb{C}P^2$.

%=============================================================================
\section{Why $S^3$ and $\mathbb{C}P^2$ Share Eigenvalues but Not
Multiplicities}
%=============================================================================

Both $S^3$ and $\mathbb{C}P^2$ have eigenvalues $\lambda_k = k(k+2)$ because:

\begin{itemize}
\item $S^3 = \mathrm{SU}(2)$ is a group manifold; its Laplacian eigenvalues
  come from the Casimir operator, giving $\lambda_k = k(k+2)$ with
  multiplicity $(k+1)^2 = (\dim V_k)^2$ where $V_k$ is the spin-$k/2$
  representation (by the Peter--Weyl theorem for $L^2(G)$).

\item $\mathbb{C}P^2 = \mathrm{SU}(3)/\mathrm{U}(2)$; its eigenvalues come
  from the Casimir of $\mathrm{SU}(3)$ restricted to spherical representations
  $(k,k)$, giving the same $\lambda_k = k(k+2)$ but with multiplicity
  $(k+1)^3 = \dim(k,k)$.
\end{itemize}

This coincidence arises because both manifolds are rank-1 symmetric spaces
whose Casimir operators yield eigenvalues of the same algebraic form $k(k+n)$
with $n=2$.  But they are geometrically distinct: $S^3$ has real dimension~3
while $\mathbb{C}P^2$ has real dimension~4.

%=============================================================================
\section{Impact on DFD Computations}
%=============================================================================

\subsection{Files Requiring Correction}

The following files use the incorrect $(k+1)^2$ multiplicity for computations
labeled as $\mathbb{C}P^2$:

\begin{enumerate}
\item \texttt{numerical\_a4.py} (lines 36, 47, 275, 300): Uses
  $(k+1)^2$ throughout.  If the intent is the $\mathbb{C}P^2$ factor of
  $X = \mathbb{C}P^2 \times S^3$, this must be changed to $(k+1)^3$.
  If the intent is the $S^3$ factor, the labeling should be corrected.

\item \texttt{a4\_Gilkey\_Computation.tex} (line 363 ff.): States
  ``multiplicity $m_l = (l+1)^2$'' and labels it as $\mathbb{C}P^2$.
  This is the $S^3$ spectrum, not $\mathbb{C}P^2$.
\end{enumerate}

\subsection{Files That Are Correct}

\begin{enumerate}
\item \texttt{Toeplitz\_Heat\_Kernel.tex}: Correctly identifies $(k+1)^3$
  for $\mathbb{C}P^2$ and $(k+1)^2$ for $S^3$, provides the
  representation-theoretic derivation, and verifies Weyl asymptotics.
\end{enumerate}

\subsection{Consequences for the Heat Kernel}

The heat kernel trace on $\mathbb{C}P^2$:
\begin{equation}
K(t) = \sum_{k=0}^{\infty} (k+1)^3\,e^{-k(k+2)t}
     = e^t \sum_{m=1}^{\infty} m^3\,e^{-m^2 t}
     \sim \frac{1}{2t^2} \quad\text{as } t\to 0^+\,,
\end{equation}
giving the correct leading Seeley--DeWitt coefficient
$a_0 = \operatorname{Vol}(\mathbb{C}P^2)/(4\pi)^2$ for a 4-dimensional
manifold.

With the incorrect $(k+1)^2$:
\begin{equation}
K_{S^3}(t) = \sum_{k=0}^{\infty} (k+1)^2\,e^{-k(k+2)t}
\sim \frac{1}{2t^{3/2}}\cdot\frac{\sqrt{\pi}}{2}
\quad\text{as }t\to 0^+\,,
\end{equation}
which gives $t^{-3/2}$ scaling, consistent with a 3-dimensional manifold.

Any $a_4$ coefficient computed with $(k+1)^2$ multiplicities but labeled
as $\mathbb{C}P^2$ will be incorrect; it will actually be the $a_4$ for $S^3$.

%=============================================================================
\section{The General Formula for $\mathbb{C}P^n$}
%=============================================================================

For completeness, the full spectrum of the scalar Laplacian on
$\mathbb{C}P^n$ with the Fubini--Study metric (normalized so that
$\mathrm{Ric} = (n{+}1)\,g$, giving $\lambda_k = 4k(k+n)$; or
$\mathrm{Ric} = 2g$ in the convention with $\lambda_k = k(k+n)$):

\begin{theorem}[Spectrum of $\Delta$ on $\mathbb{C}P^n$]
\label{thm:CPn-spectrum}
The eigenvalues are $\lambda_k = k(k+n)$ for $k = 0, 1, 2, \ldots$
with multiplicities
\begin{equation}
m_k = \binom{k+n}{n}^{\!2} - \binom{k+n-1}{n}^{\!2}
    = \frac{2k+n}{n}\,\binom{k+n-1}{n-1}^{\!2}\,.
\end{equation}
\end{theorem}

\noindent
This follows from the decomposition of $L^2(\mathbb{C}P^n)$ into spherical
representations of $\mathrm{SU}(n{+}1)$ with respect to $\mathrm{U}(n)$,
combined with the Weyl dimension formula for $\mathrm{SU}(n{+}1)$.

\paragraph{References.}
\begin{itemize}
\item Berger, Gauduchon, Mazet, \textit{Le Spectre d'une Vari\'et\'e
  Riemannienne}, Lecture Notes in Mathematics 194, Springer, 1971.
\item Ikeda, Taniguchi, ``Spectra and eigenforms of the Laplacian on
  $S^n$ and $P^n(\mathbb{C})$,'' \textit{Osaka J.\ Math.}\ \textbf{15}
  (1978), 515--546.
\item Lu Qikeng, ``The eigenfunctions of the complex projective space,''
  \textit{Acta Mathematica Sinica, English Series}, 1998.
\end{itemize}

%=============================================================================
\section{Table of First Few Eigenvalues and Multiplicities}
%=============================================================================

\begin{center}
\begin{tabular}{crrrrl}
\toprule
$k$ & $\lambda_k = k(k+2)$ & $m_k^{S^3} = (k+1)^2$ & $m_k^{\mathbb{C}P^2}
= (k+1)^3$ & Cumulative $N$ & Comment \\
\midrule
0 &   0 &   1 &    1 &      1 & Zero mode \\
1 &   3 &   4 &    8 &      9 & \\
2 &   8 &   9 &   27 &     36 & \\
3 &  15 &  16 &   64 &    100 & \\
4 &  24 &  25 &  125 &    225 & \\
5 &  35 &  36 &  216 &    441 & \\
6 &  48 &  49 &  343 &    784 & \\
7 &  63 &  64 &  512 &  1,296 & \\
8 &  80 &  81 &  729 &  2,025 & \\
9 &  99 & 100 & 1000 &  3,025 & \\
\bottomrule
\end{tabular}
\end{center}

Note that the cumulative count with $(k+1)^3$ is
$N = \sum_{j=0}^{k}(j+1)^3 = \bigl[\frac{(k+1)(k+2)}{2}\bigr]^2
= \binom{k+2}{2}^2$, a perfect square---this is the total dimension of
bi-degree $\leq k$ polynomial space, as expected.

\end{document}
